\par\clearpage\SourceLeaf{323}\ArithmeticHead{14}{L}{A. OSTROWSKI, ÜBER DEN ERSTEN UND VIERTEN GAUSSSCHEN BEWEIS USW.}
\noindent\hspace*{\parindent}\textls[100]{4. Zusammenfassung der Punkte von \(U=0\) innerhalb \((R)\) zu
Elementarbogen.} Für ein \(x\) mit \(x_1<x<x_2\) bezeichnen wir die innerhalb
\((R)\) liegende Wurzel von \(U^{(0)}\) mit der kleinsten Ordinate durch \(W_1(x)\),
diejenige mit der zweitkleinsten Ordinate durch \(W_2(x)\) usw. So erhalten wir
\(m\) eindeutige Funktionen \(W_1(x),\ W_2(x),\ \ldots,\ W_m(x)\), von denen nun aber
sofort folgt, dass sie für \(x_1<x<x_2\) stetig sind. Um dies etwa für \(W_i(x)\) zu
beweisen, brauchen wir nur zu zeigen, dass für jedes \(\overline{x}\) zwischen \(x_1\) und \(x_2\)
sich jedem hinreichend klein vorgegebenen \(\varepsilon>0\) ein solches \(\delta\) zuordnen lässt,
dass die der Grösse nach \(i\)-te Wurzel von \(U^{(0)}\) innerhalb \((R)\), die \(x\) entspricht,
sich von \(W_i(\overline{x})\) um weniger als \(\varepsilon\) unterscheidet, solange \(\overline{x}-\delta<x<\overline{x}+\delta\)
bleibt. Dies haben wir aber in Nr. 3 nachgewiesen. Jetzt folgt weiter aus
dem bekannten Satz der Differentialrechnung, dass \(W_i(x)\) differenzierbar und
dass
\[
\frac{\partial W_i(x)}{\partial x}=-\frac{\dfrac{\partial U^{(0)}}{\partial x}}{\dfrac{\partial U^{(0)}}{\partial y}}
\]
ist, also zwischen \(x_1\) und \(x_2\) dasselbe Vorzeichen behält. Folglich sind die \(W_i\)
monoton. Nähert sich \(x\) etwa dem Werte \(x_1\), so strebt \(W_i(x)\) einem ganz
bestimmten Grenzwerte zu, da ja \(|W_i(x)|\leqq R\) gilt. —
\par
Genau die gleichen Überlegungen gelten auch für jedes der Intervalle\linebreak
\((x_2,x_3),\ldots\,(x_{k-1},x_k)\) und ebenso auch für die beiden Kreissegmente, die ev.
nach links von \(x_1\), bezw. nach rechts von \(x_k\) liegen. Damit ist die ganze
Kurve \(U^{(0)}=0\) innerhalb von \((R)\) in eine endliche Anzahl regulärer Kurvenbogen
zerlegt (wir werden sie gelegentlich als \textls[100]{Elementarbogen} bezeichnen),
zu denen noch isolierte Punkte auf den Geraden \(x=x_1,\ x=x_2,\ \ldots\) in endlicher
Anzahl (höchstens \(n\) auf jeder) hinzukommen können. (Man kann zwar
zeigen, dass solche isolierte Punkte nicht vorkommen, doch dies ist für uns
unwesentlich.) — Unser Resultat gilt in dieser Formulierung auch für \(U=0\)
(sowie für \(U^{(1)}=0\)), da die eventuell hinzukommenden Geraden \(x=x'\) durch
ihre Schnittpunkte mit den \(W_i\)-Kurven in je endlich viele Elementarbogen
zerlegt werden. —
\par
Innerhalb des von den Parallelen \(x=x_1,\ x=x_2\) und den Elementarbogen
\(y=W_i(x),\ y=W_{i+1}(x)\) begrenzten Bereiches (in einigen Fällen muss einer
oder beide Bogen durch die Bogen des Kreises \((R)\) ersetzt werden) hat \(U^{(0)}\)\par
\clearpage\SourceLeaf{324}\ArithmeticHead{15}{R}{A. OSTROWSKI, ÜBER DEN ERSTEN UND VIERTEN GAUSSSCHEN BEWEIS USW.}
\noindent stets dasselbe Vorzeichen, und dieses Vorzeichen wechselt längs eines Bogens
\(y=W_i(x)\), da auf ihm \(\frac{\partial U^{(0)}}{\partial x}\ne0\) ist.
\par
Daraus folgt ferner, dass auch \(U^{(1)}\) sein Vorzeichen längs der Elementarbogen
wechselt, in die die Kurve \(U=0\) innerhalb \((R)\) zerfällt.
\par
Es sei nun \((x_1,w)\) ein Punkt auf der Geraden \(x=x_1\), der nicht auf dem
Kreise \((R)\) liegt und in dem von links und von rechts mehrere Elementarbogen
zusammenkommen, etwa
\[
y=W_i(x),\ \ldots,\ y=W_{i+k}(x)
\]
von rechts und
\[
y=\overline{W}_{i'}(x),\ y=\overline{W}_{i'+1}(x),\ \ldots,\ y=\overline{W}_{i'+k'}(x)
\]
von links. Wir grenzen auf \(x=x_1\) um \((x_1,w)\) zwei Intervalle \(J_+(w+\varepsilon>y>w)\)
und \(J_-(w-\varepsilon<y<w)\) ab, in denen keine weitere Wurzel von \(U\) liegt. Dann
grenzt längst \(J_+\) das Gebiet zwischen \(y=W_{i+k}(x)\) und \(y=W_{i+k+1}(x)\) an
das Gebiet zwischen \(y=\overline{W}_{i'+k'}(x)\) und \(y=\overline{W}_{i'+k'+1}(x)\), und das Vorzeichen
von \(U^{(0)}\) ist in den beiden Gebieten dasselbe, wie in \(J_+\). Diese beiden Gebiete
bilden daher ein Gebiet konstanten Vorzeichens von \(U^{(0)}\). Dasselbe gilt auch
für die Gebiete, die längst \(J_-\) aneinander grenzen. Da daher \(U^{(0)}\) in den an
den Punkt \((x_1,w)\) anstossenden Gebieten abwechselndes Vorzeichen hat, so ist
die Anzahl dieser Gebiete \textls[100]{gerade}. Daher ist auch die Anzahl der in diesem
Punkt zusammenkommenden Bogen \textls[100]{gerade}. Und diese Anzahl bleibt auch
dann gerade, wenn man nachträglich die durch \((x_1,w)\) gehende Gerade \(x=x_1\)
hinzufügen muss, d. h. für \(U=0\). Was aber die auf dem Kreise \((R)\) liegenden
Punkte von \(U^{(1)}=0\) betrifft, so können wir annehmen, dass sich in ihnen an
jeden ausserhalb des Kreises verlaufenden Zweig von \(U=0\) genau ein Elementarbogen
anschliesst. Denn dies ist nach dem eingangs über diese Zweige
Bewiesenen stets zu erreichen, indem man \(R\) etwas grösser annimmt.
\par
